\documentclass[conference]{IEEEtran}

\usepackage{amsmath,amssymb}
\usepackage{newtxtext,newtxmath}
\usepackage{booktabs}
\usepackage{multirow}
\usepackage{graphicx}
\usepackage{microtype}
\usepackage{xcolor}
\usepackage[hidelinks]{hyperref}
\usepackage{url}
\usepackage{balance}

\graphicspath{{./}{../results/figures/}}
% Central source of every reported scalar in the REIM manuscript.
%
% Values in "completed component diagnostics" are read from immutable JSON
% artifacts in results/tables and datasets.  Final benchmark values are filled
% only by scripts/update_paper_results.py after its raw-episode publication
% gate succeeds.

\newif\ifREIMFinalResults
\REIMFinalResultstrue

% MT10/MT50 paper claims have an independent publication gate.  It remains
% closed until both official-clean suites and every disturbed extension pass
% completeness, provenance, and paired-protocol validation.
\newif\ifREIMMultiTaskResults
\REIMMultiTaskResultsfalse

\newcommand{\PendingMetric}{\textnormal{\textsc{tbd}}}
\newcommand{\PendingInterval}{[\textnormal{\textsc{tbd}}]}
\newcommand{\NotApplicable}{--}

% ---------------------------------------------------------------------------
% Verified environment, data, and model configuration
% ---------------------------------------------------------------------------
\newcommand{\REIMTask}{\texttt{pick-place-v3}}
\newcommand{\REIMStateDim}{21}
\newcommand{\REIMActionDim}{4}
\newcommand{\REIMDemoTrajectories}{500}
\newcommand{\REIMDemoTransitions}{25{,}969}
\newcommand{\REIMDemoSuccess}{100\%}
\newcommand{\REIMDevelopmentSeed}{42}
\newcommand{\DevelopmentTrainingSeeds}{3}
\newcommand{\ConfidenceLevel}{95\%}

\newcommand{\ACTChunkLength}{20}
\newcommand{\ACTHiddenDim}{256}
\newcommand{\ACTAttentionHeads}{8}
\newcommand{\ACTEncoderLayers}{3}
\newcommand{\ACTDecoderLayers}{4}
\newcommand{\ACTFeedForwardDim}{1{,}024}
\newcommand{\ACTLatentDim}{32}
\newcommand{\ACTTrainingEpochs}{200}
\newcommand{\ACTParameterCount}{6{,}627{,}908}

\newcommand{\DetectorHistory}{10}
\newcommand{\DetectorHiddenDim}{128}
\newcommand{\DetectorMLPDim}{64}
\newcommand{\DetectorTrainingEpisodes}{2{,}000}
\newcommand{\DetectorTrainingSamples}{285{,}443}
\newcommand{\DetectorValidationSamples}{56{,}253}
\newcommand{\DetectorForecastHorizon}{10}
\newcommand{\DetectorSelectionThreshold}{0.50}
\newcommand{\DetectorDiagnosticThreshold}{0.20}

\newcommand{\CurriculumTriggerThreshold}{0.10}
\newcommand{\CurriculumNoiseLevel}{20\%}
\newcommand{\CurriculumActionNoise}{0.08}
\newcommand{\CurriculumObservationNoise}{0.005}
\newcommand{\CurriculumObjectImpulse}{0.02}
\newcommand{\RecoveryTrainEpisodes}{1{,}000}
\newcommand{\RecoveryValidationEpisodes}{200}
\newcommand{\RecoveryTrainStarts}{985}
\newcommand{\RecoveryValidationStarts}{193}
\newcommand{\RecoveryTrainEligibleStarts}{978}
\newcommand{\RecoveryValidationEligibleStarts}{192}
\newcommand{\RecoveryTrainExpertSamples}{42{,}386}
\newcommand{\RecoveryValidationExpertSamples}{8{,}212}
\newcommand{\RecoveryExpertEpochs}{40}

\newcommand{\PPOTrainingTarget}{500{,}000}
\newcommand{\PPODevelopmentSteps}{500{,}136}
\newcommand{\PPOSelectedStep}{118{,}248}
\newcommand{\PPOEvalEpisodes}{50}
\newcommand{\PPORecoveryBudget}{150}
\newcommand{\PPOHiddenDim}{256}
\newcommand{\PPOParallelEnvs}{4}
\newcommand{\PPORolloutSteps}{1{,}024}
\newcommand{\PPOBatchSize}{512}
\newcommand{\PPOEpochsPerUpdate}{5}
\newcommand{\PPOLearningRate}{3\times10^{-5}}
\newcommand{\PPOGamma}{0.99}
\newcommand{\PPOGAE}{0.95}
\newcommand{\PPOClip}{0.10}
\newcommand{\PPOEntropyCoefficient}{0.0001}

\newcommand{\TaskSuccessBonus}{25}
\newcommand{\TerminalFailurePenalty}{-10}
\newcommand{\TimePenalty}{-0.01}
\newcommand{\BaseRewardScale}{1}
\newcommand{\GoalProgressScale}{2}
\newcommand{\ReachProgressScale}{2}
\newcommand{\LiftProgressScale}{2}

% Frozen end-to-end controller.  These values are deliberately distinct from
% the lower threshold used above to collect a broad early-trigger curriculum.
% min_steps == budget and clear_steps > budget disable any intermediate
% detector-based release: the recovery option holds control until task success
% or budget expiry.
\newcommand{\FinalControllerTrigger}{0.20}
\newcommand{\FinalControllerRelease}{0}
\newcommand{\FinalControllerMinSteps}{150}
\newcommand{\FinalControllerBudget}{150}
\newcommand{\FinalControllerClearSteps}{200}

\newcommand{\PlannedMainEpisodes}{1{,}000}
\newcommand{\PlannedRobustnessEpisodes}{200}
\newcommand{\RobustnessLevels}{0\%, 10\%, 20\%, 30\%, and 40\%}

% Shared known-task MT10/MT50 protocol.  These are configuration constants,
% not experimental outcomes.
\newcommand{\MTTenBenchmark}{\textsc{MT10}}
\newcommand{\MTFiftyBenchmark}{\textsc{MT50}}
\newcommand{\MTRawObservationDim}{39}
\newcommand{\MTTenStateDim}{49}
\newcommand{\MTFiftyStateDim}{89}
\newcommand{\MTPerTaskDemonstrations}{50}
\newcommand{\MTPerTaskEvaluationEpisodes}{50}
\newcommand{\MTMaxEpisodeSteps}{500}
\newcommand{\MTDisturbanceLevels}{0\%, 10\%, 20\%, 30\%, and 40\%}
\newcommand{\MTFailureThresholdQuantile}{0.90}
\newcommand{\MTFailureForecastHorizon}{10}
\newcommand{\MTTerminalFailureTail}{25}

% ---------------------------------------------------------------------------
% Completed component diagnostics (not final end-to-end benchmark evidence)
% ---------------------------------------------------------------------------
\newcommand{\ACTBestValidationLone}{0.004013}
\newcommand{\ACTBestValidationLoss}{0.004047}

\newcommand{\DetectorAccuracyHalf}{89.3\%}
\newcommand{\DetectorPrecisionHalf}{96.5\%}
\newcommand{\DetectorRecallHalf}{86.8\%}
\newcommand{\DetectorFoneHalf}{91.4\%}
\newcommand{\DetectorAccuracyDiagnostic}{89.6\%}
\newcommand{\DetectorPrecisionDiagnostic}{91.8\%}
\newcommand{\DetectorRecallDiagnostic}{92.3\%}
\newcommand{\DetectorFoneDiagnostic}{92.1\%}
\newcommand{\DetectorTNDiagnostic}{16{,}352}
\newcommand{\DetectorFPDiagnostic}{3{,}038}
\newcommand{\DetectorFNDiagnostic}{2{,}832}
\newcommand{\DetectorTPDiagnostic}{34{,}031}
\newcommand{\DetectorCurrentEventPositiveShare}{86.1\%}
\newcommand{\DetectorStrictPreEventPrecision}{53.5\%}
\newcommand{\DetectorStrictPreEventRecall}{68.3\%}
\newcommand{\DetectorStrictPreEventFone}{60.0\%}
\newcommand{\DetectorPreEventDetectedTrajectories}{191}
\newcommand{\DetectorEventTrajectories}{258}
\newcommand{\DetectorPreEventTrajectoryDetection}{74.0\%}
\newcommand{\DetectorMedianPreEventLead}{1}
\newcommand{\DetectorNonEventTrajectoryAlert}{66.2\%}

\newcommand{\RecoveryWarmStartValidationLoss}{0.006391}
\newcommand{\RecoverySelectionSuccess}{100.0\%}
\newcommand{\RecoverySelectionRecovery}{100.0\%}
\newcommand{\RecoverySelectionFailure}{0.0\%}

% Provenance for the qualitative paired simulator trace.  This is not an
% aggregate benchmark result and uses its own explicitly reported gate.
\newcommand{\OperationTraceSeed}{8{,}300{,}042}
\newcommand{\OperationTraceNoise}{20.0\%}
\newcommand{\OperationTraceThreshold}{0.20}
\newcommand{\OperationTraceACTSteps}{200}
\newcommand{\OperationTraceREIMSteps}{62}
\newcommand{\OperationTraceRecoverySteps}{53}

% ---------------------------------------------------------------------------
% Final paired end-to-end benchmark results. These commands are populated by
% scripts/update_paper_results.py after validating every canonical raw record.
% Success and recovery values are percentages. For learned recovery methods,
% "Recovery" is task completion while recovery has control, divided by
% intervention count.
% Random Reset reports post-reset completion per reset and is labeled separately.
% ---------------------------------------------------------------------------
\newcommand{\FinalACTSuccess}{73.4\%}
\newcommand{\FinalACTSuccessCI}{[70.6, 76.0]\%}
\newcommand{\FinalACTRecovery}{\NotApplicable}
\newcommand{\FinalACTRecoveryCI}{}
\newcommand{\FinalACTIntervention}{\NotApplicable}
\newcommand{\FinalACTSteps}{93.3}

\newcommand{\FinalResetSuccess}{81.3\%}
\newcommand{\FinalResetSuccessCI}{[78.8, 83.6]\%}
\newcommand{\FinalResetRecovery}{74.5\%}
\newcommand{\FinalResetRecoveryCI}{[71.3, 77.6]\%}
\newcommand{\FinalResetIntervention}{73.2\%}
\newcommand{\FinalResetSteps}{93.3}

\newcommand{\FinalRLRecoverySuccess}{87.7\%}
\newcommand{\FinalRLRecoverySuccessCI}{[85.5, 89.6]\%}
\newcommand{\FinalRLRecoveryRate}{83.3\%}
\newcommand{\FinalRLRecoveryRateCI}{[80.6, 85.9]\%}
\newcommand{\FinalRLRecoveryIntervention}{73.1\%}
\newcommand{\FinalRLRecoverySteps}{72.8}

\newcommand{\FinalREIMSuccess}{90.4\%}
\newcommand{\FinalREIMSuccessCI}{[88.4, 92.1]\%}
\newcommand{\FinalREIMRecovery}{85.3\%}
\newcommand{\FinalREIMRecoveryCI}{[82.5, 88.0]\%}
\newcommand{\FinalREIMIntervention}{65.5\%}
\newcommand{\FinalREIMSteps}{69.3}
\newcommand{\FinalREIMGain}{+17.0\,pp}
\newcommand{\FinalREIMGainCI}{[+14.7, +19.3]\,pp}
\newcommand{\FinalREIMRescuedEpisodes}{170}
\newcommand{\FinalREIMHarmedEpisodes}{0}
\newcommand{\FinalREIMInterventionDeltaVsRL}{-7.6\,pp}
\newcommand{\FinalREIMInterventionReductionVsRL}{10.4\%}
\newcommand{\FinalREIMFewerIntervenedEpisodes}{76}
\newcommand{\FinalREIMGainVsReset}{+9.1\,pp}
\newcommand{\FinalREIMGainVsResetCI}{[+6.1, +12.1]\,pp}
\newcommand{\FinalREIMDeltaVsRL}{+2.7\,pp}
\newcommand{\FinalREIMDeltaVsRLCI}{[+0.7, +4.8]\,pp}

% Final component ablation.
\newcommand{\AblationASuccess}{73.4\%}
\newcommand{\AblationARecovery}{\NotApplicable}
\newcommand{\AblationASteps}{93.3}
\newcommand{\AblationBSuccess}{49.1\%}
\newcommand{\AblationBRecovery}{\NotApplicable}
\newcommand{\AblationBSteps}{129.5}
\newcommand{\AblationCSuccess}{87.7\%}
\newcommand{\AblationCRecovery}{83.3\%}
\newcommand{\AblationCSteps}{72.8}
\newcommand{\AblationDSuccess}{90.4\%}
\newcommand{\AblationDRecovery}{85.3\%}
\newcommand{\AblationDSteps}{69.3}

% Final robustness endpoints used in the textual summary.  The complete curve
% remains an external figure input so that it is regenerated from episode CSV.
\newcommand{\FinalACTSuccessAtZeroNoise}{100.0\%}
\newcommand{\FinalACTSuccessAtFortyNoise}{31.5\%}
\newcommand{\FinalREIMSuccessAtZeroNoise}{100.0\%}
\newcommand{\FinalREIMSuccessAtFortyNoise}{63.5\%}
\newcommand{\FinalREIMGainAtTenNoise}{+2.5\,pp}
\newcommand{\FinalREIMGainAtTwentyNoise}{+15.5\,pp}
\newcommand{\FinalREIMGainAtThirtyNoise}{+26.0\,pp}
\newcommand{\FinalREIMGainAtFortyNoise}{+32.0\,pp}

% ---------------------------------------------------------------------------
% Pending MT10/MT50 breadth results.  These placeholders are never rendered
% while \ifREIMMultiTaskResults is false.  A future publication-gate script
% must populate every value atomically before enabling the gate.
% ---------------------------------------------------------------------------
\newcommand{\MTTenMLPCleanSuccess}{\PendingMetric}
\newcommand{\MTTenMLPCleanSuccessCI}{\PendingInterval}
\newcommand{\MTTenACTCleanSuccess}{\PendingMetric}
\newcommand{\MTTenACTCleanSuccessCI}{\PendingInterval}
\newcommand{\MTTenHeuristicCleanSuccess}{\PendingMetric}
\newcommand{\MTTenHeuristicCleanSuccessCI}{\PendingInterval}
\newcommand{\MTTenREIMCleanSuccess}{\PendingMetric}
\newcommand{\MTTenREIMCleanSuccessCI}{\PendingInterval}
\newcommand{\MTFiftyMLPCleanSuccess}{\PendingMetric}
\newcommand{\MTFiftyMLPCleanSuccessCI}{\PendingInterval}
\newcommand{\MTFiftyACTCleanSuccess}{\PendingMetric}
\newcommand{\MTFiftyACTCleanSuccessCI}{\PendingInterval}
\newcommand{\MTFiftyHeuristicCleanSuccess}{\PendingMetric}
\newcommand{\MTFiftyHeuristicCleanSuccessCI}{\PendingInterval}
\newcommand{\MTFiftyREIMCleanSuccess}{\PendingMetric}
\newcommand{\MTFiftyREIMCleanSuccessCI}{\PendingInterval}

\newcommand{\MTTenMLPWorstQuartile}{\PendingMetric}
\newcommand{\MTTenACTWorstQuartile}{\PendingMetric}
\newcommand{\MTTenHeuristicWorstQuartile}{\PendingMetric}
\newcommand{\MTTenREIMWorstQuartile}{\PendingMetric}
\newcommand{\MTFiftyMLPWorstQuartile}{\PendingMetric}
\newcommand{\MTFiftyACTWorstQuartile}{\PendingMetric}
\newcommand{\MTFiftyHeuristicWorstQuartile}{\PendingMetric}
\newcommand{\MTFiftyREIMWorstQuartile}{\PendingMetric}

\newcommand{\MTTenHeuristicIntervention}{\PendingMetric}
\newcommand{\MTTenREIMIntervention}{\PendingMetric}
\newcommand{\MTFiftyHeuristicIntervention}{\PendingMetric}
\newcommand{\MTFiftyREIMIntervention}{\PendingMetric}

\newcommand{\MTTenREIMDeltaVsACT}{\PendingMetric}
\newcommand{\MTTenREIMDeltaVsACTCI}{\PendingInterval}
\newcommand{\MTTenMLPDeltaVsACT}{\PendingMetric}
\newcommand{\MTTenMLPDeltaVsACTCI}{\PendingInterval}
\newcommand{\MTTenHeuristicDeltaVsACT}{\PendingMetric}
\newcommand{\MTTenHeuristicDeltaVsACTCI}{\PendingInterval}
\newcommand{\MTFiftyREIMDeltaVsACT}{\PendingMetric}
\newcommand{\MTFiftyREIMDeltaVsACTCI}{\PendingInterval}
\newcommand{\MTFiftyMLPDeltaVsACT}{\PendingMetric}
\newcommand{\MTFiftyMLPDeltaVsACTCI}{\PendingInterval}
\newcommand{\MTFiftyHeuristicDeltaVsACT}{\PendingMetric}
\newcommand{\MTFiftyHeuristicDeltaVsACTCI}{\PendingInterval}

% This generated file is absent until the independent publication gate has
% validated both suites, all disturbance conditions, and both five-bank
% separation audits.  A fresh or incomplete run therefore remains closed.
\InputIfFileExists{multitask_results.tex}{% Generated by visualization/generate_multitask_paper_assets.py.
% The publication gate validated MT10 + MT50 clean, five robustness
% conditions per suite, and both independent five-bank audits.
% input-manifest-sha256: 0a5063866bc86473be04d8fb3f36f0201422de8a283caf3997ba487dd70b4e4d
\REIMMultiTaskResultstrue
\renewcommand{\MTTenMLPCleanSuccess}{94.2\%}
\renewcommand{\MTTenMLPCleanSuccessCI}{[92.4, 96.0]\%}
\renewcommand{\MTTenMLPWorstQuartile}{80.7\%}
\renewcommand{\MTTenMLPDeltaVsACT}{-3.2\,pp}
\renewcommand{\MTTenMLPDeltaVsACTCI}{[-5.2, -1.4]\,pp}
\renewcommand{\MTTenACTCleanSuccess}{97.4\%}
\renewcommand{\MTTenACTCleanSuccessCI}{[96.0, 98.6]\%}
\renewcommand{\MTTenACTWorstQuartile}{91.3\%}
\renewcommand{\MTTenHeuristicCleanSuccess}{98.0\%}
\renewcommand{\MTTenHeuristicCleanSuccessCI}{[96.8, 99.0]\%}
\renewcommand{\MTTenHeuristicWorstQuartile}{93.3\%}
\renewcommand{\MTTenHeuristicIntervention}{23.4\%}
\renewcommand{\MTTenHeuristicOccupancy}{14.2\%}
\renewcommand{\MTTenHeuristicDeltaVsACT}{+0.6\,pp}
\renewcommand{\MTTenHeuristicDeltaVsACTCI}{[-1.0, +2.2]\,pp}
\renewcommand{\MTTenREIMCleanSuccess}{97.4\%}
\renewcommand{\MTTenREIMCleanSuccessCI}{[96.0, 98.6]\%}
\renewcommand{\MTTenREIMWorstQuartile}{91.3\%}
\renewcommand{\MTTenREIMIntervention}{36.4\%}
\renewcommand{\MTTenREIMOccupancy}{14.3\%}
\renewcommand{\MTTenREIMDeltaVsACT}{+0.0\,pp}
\renewcommand{\MTTenREIMDeltaVsACTCI}{[+0.0, +0.0]\,pp}
\renewcommand{\MTFiftyMLPCleanSuccess}{81.5\%}
\renewcommand{\MTFiftyMLPCleanSuccessCI}{[80.4, 82.6]\%}
\renewcommand{\MTFiftyMLPWorstQuartile}{41.5\%}
\renewcommand{\MTFiftyMLPDeltaVsACT}{-10.2\,pp}
\renewcommand{\MTFiftyMLPDeltaVsACTCI}{[-11.4, -8.9]\,pp}
\renewcommand{\MTFiftyACTCleanSuccess}{91.6\%}
\renewcommand{\MTFiftyACTCleanSuccessCI}{[90.8, 92.4]\%}
\renewcommand{\MTFiftyACTWorstQuartile}{69.1\%}
\renewcommand{\MTFiftyHeuristicCleanSuccess}{94.0\%}
\renewcommand{\MTFiftyHeuristicCleanSuccessCI}{[93.3, 94.7]\%}
\renewcommand{\MTFiftyHeuristicWorstQuartile}{77.1\%}
\renewcommand{\MTFiftyHeuristicIntervention}{17.8\%}
\renewcommand{\MTFiftyHeuristicOccupancy}{10.6\%}
\renewcommand{\MTFiftyHeuristicDeltaVsACT}{+2.4\,pp}
\renewcommand{\MTFiftyHeuristicDeltaVsACTCI}{[+1.6, +3.0]\,pp}
\renewcommand{\MTFiftyREIMCleanSuccess}{92.3\%}
\renewcommand{\MTFiftyREIMCleanSuccessCI}{[91.4, 93.1]\%}
\renewcommand{\MTFiftyREIMWorstQuartile}{71.7\%}
\renewcommand{\MTFiftyREIMIntervention}{20.2\%}
\renewcommand{\MTFiftyREIMOccupancy}{11.0\%}
\renewcommand{\MTFiftyREIMDeltaVsACT}{+0.6\,pp}
\renewcommand{\MTFiftyREIMDeltaVsACTCI}{[+0.2, +1.1]\,pp}
}{}

% Separate post-freeze gate audit. The matched heuristic point uses the same
% serialized 200-episode CRN bank as the LSTM-threshold sweep.
\newcommand{\GateHeuristicSuccess}{86.0\%}
\newcommand{\GateHeuristicIntervention}{76.0\%}
\newcommand{\GateTauSeventeenFiveSuccess}{92.5\%}
\newcommand{\GateTauSeventeenFiveIntervention}{76.5\%}
\newcommand{\GateTauTwentySuccess}{90.5\%}
\newcommand{\GateTauTwentyIntervention}{63.5\%}
\newcommand{\GateTauSeventeenFiveGainVsHeuristic}{+6.5\,pp}
\newcommand{\GateTauSeventeenFiveGainVsHeuristicCI}{[+2.5, +11.0]\,pp}
\newcommand{\GateTauSeventeenFiveP}{0.0044}
\newcommand{\GateTauTwentyGainVsHeuristic}{+4.5\,pp}
\newcommand{\GateTauTwentyGainVsHeuristicCI}{[+0.5, +8.5]\,pp}
\newcommand{\GateTauTwentyP}{0.0490}
\newcommand{\GateTauSeventeenFiveInterventionDelta}{+0.5\,pp}
\newcommand{\GateTauSeventeenFiveInterventionDeltaCI}{[-8.0, +8.5]\,pp}


\newcommand{\policyact}{\pi_{\mathrm{ACT}}}
\newcommand{\policyrec}{\pi_{\mathrm{rec}}}
\newcommand{\PaperGraphic}[3]{%
  \IfFileExists{#1.pdf}{%
    \includegraphics[width=#2]{#1.pdf}%
  }{%
    \IfFileExists{#1.png}{%
      \includegraphics[width=#2]{#1.png}%
    }{%
      \fbox{\parbox[c][1.15in][c]{0.92\linewidth}{%
        \centering\textbf{Missing compiled artifact}\\[2pt]#3}}%
    }%
  }%
}

\title{REIM: Recovery-Enhanced Imitation Learning for\\
Failure-Aware Manipulation in Meta-World}

\author{\IEEEauthorblockN{Anonymous Authors}
\IEEEauthorblockA{Submission under review}}

\begin{document}
\maketitle

\begin{abstract}
Imitation-learned manipulation policies can be accurate on demonstrations yet
brittle after a missed grasp, object drop, or control perturbation shifts the
closed-loop state distribution. We introduce Recovery-Enhanced Imitation
Learning (REIM), a failure-aware controller that combines a state-based
Action-Chunking-with-Transformers (ACT) task policy, a causal LSTM risk monitor,
and a trigger-aligned supervised recovery option. Its curriculum executes ACT
online under disturbances, serializes the complete MuJoCo simulator state at an
early collection trigger, and queries a scripted expert for successful
continuations from those learner-induced states. A deterministic recovery MLP
is trained with Smooth L1 loss on 42,386 state--action pairs and selected on
8,212 disjoint validation pairs. About 99\% of the captured states precede any
lift, so REIM is principally an early pre-grasp correction system rather than a
post-drop repair policy.
Across \PlannedMainEpisodes{} paired Meta-World PickPlace trials at 20\%
mixed simulated disturbance, REIM reaches \FinalREIMSuccess{} task success
versus \FinalACTSuccess{} for ACT and \FinalResetSuccess{} for one fresh-task
retry. The same recovery actor with a semantic heuristic gate reaches
\FinalRLRecoverySuccess{}, while REIM improves success by
\FinalREIMDeltaVsRL{} and uses
\FinalREIMInterventionReductionVsRL{} fewer interventions. At 40\%
disturbance, REIM retains
\FinalREIMSuccessAtFortyNoise{} success versus
\FinalACTSuccessAtFortyNoise{} for ACT.
We additionally define a shared task-conditioned instantiation for the
goal-observable, known-task Meta-World MT10 and MT50 suites, with official-clean
evaluation separated from a task-universal disturbance extension.
All experiments in this paper use the simulated Sawyer model in
Meta-World/MuJoCo; no physical-robot evaluation is claimed.
\end{abstract}

\section{Introduction}
\label{sec:introduction}

Imitation learning offers a direct route from expert demonstrations to robot
control, but supervised action error does not capture the distribution induced
by the learned policy. Small errors alter future observations and compound over
the horizon, a problem formalized by interactive imitation-learning work such
as DAgger~\cite{ross2011dagger}. Action chunking reduces the effective horizon
and has enabled precise manipulation~\cite{zhao2023act}, but a chunked policy
can still enter states absent from the successful demonstrations. In PickPlace,
typical departures include a failed enclosure, premature gripper opening,
object displacement, and trajectories that cease to make goal progress.

A useful recovery system must answer three coupled questions: when nominal
execution has become risky, which controller should act next, and which states
should train that controller. Training a recovery policy from hand-designed
``failure-like'' resets leaves a train--deploy gap: the reset may match object
position while omitting robot velocity, gripper dynamics, controller memory,
or task bookkeeping. Conversely, corrective demonstrations collected only
from nominal states do not teach a policy how to respond at the learner-induced
states where intervention is actually required.

REIM addresses these issues with a closed-loop composition of imitation,
runtime risk monitoring, and targeted recovery imitation. Its main
contributions are:

\begin{itemize}
  \item a causal risk-gated controller coupling temporally ensembled ACT, an
  LSTM risk monitor, and a persistent deterministic recovery option;
  \item an online trigger-state curriculum that serializes and restores full
  numeric Meta-World/MuJoCo states observed when disturbed ACT first crosses a
  broad risk-monitor collection threshold;
  \item trigger-aligned supervised recovery trained on successful expert
  continuations with a trajectory-disjoint validation split, exported as a
  standalone actor with zero reinforcement-learning interactions; and
  \item a two-level reproducible evaluation design: a PickPlace mechanism study
  separating classifier, recovery, success, and intervention evidence, plus a
  publication-gated MT10/MT50 known-task breadth protocol that keeps official
  clean scores distinct from task-universal disturbed evaluation.
\end{itemize}

\section{Related Work}
\label{sec:related}

\subsection{Imitation Learning Under Distribution Shift}

Behavior cloning is vulnerable to covariate shift because its training states
come from the expert while its test states are induced by the learner.
Interactive aggregation methods query corrective actions on learner-induced
states~\cite{ross2011dagger}. ACT instead predicts temporally coherent action
chunks with a conditional variational Transformer and ensembles overlapping
predictions~\cite{zhao2023act}. REIM retains ACT as a strong nominal policy and
adds an explicit recovery route for states that successful demonstrations do
not cover.

\subsection{Runtime Failure Monitoring}

Runtime monitors can use predictive uncertainty, action consistency, semantic
progress, or supervised failure labels. Recent work separates erratic actions
from stalled task progress~\cite{agia2025sentinel}. REIM uses a lightweight
causal LSTM~\cite{hochreiter1997lstm} over semantic state history. It estimates
the risk that a labeled event is present now or occurs within a short horizon;
it never consumes future observations at deployment.

\subsection{Recovery Policies}

Recovery RL separates a task policy from a safety-oriented recovery
policy~\cite{thananjeyan2021recovery}. Residual RL adds learned corrective
control to a nominal controller~\cite{johannink2019residual}, while
MetaReSkill~\cite{vats2023metareskill} allocates training to manipulation
failure modes and RecoveryChaining~\cite{vats2024recoverychaining} learns local
policies that reconnect to nominal skills. REIM differs in the source of its
recovery initial-state distribution: each start is a complete simulator state
produced by the ACT--detector pair at its first online collection trigger. The
collection gate is deliberately broader than the frozen deployment gate. This
permits targeted corrective imitation from learner-induced states without
requiring on-policy recovery interactions. State restoration is used only to
construct a controlled simulation dataset; we do not claim that the same reset
mechanism is directly available on hardware.

\section{Method}
\label{sec:method}

\subsection{Problem Formulation}

We consider a finite-horizon manipulation MDP with semantic state
$s_t\in\mathbb{R}^{\REIMStateDim}$ and Cartesian action
$a_t\in[-1,1]^{\REIMActionDim}$. The state concatenates the simulated Sawyer
joint positions, end-effector position and quaternion, object position, goal
position, and gripper state. The action contains end-effector displacement
$(\Delta x,\Delta y,\Delta z)$ and gripper command $g$. Task success is the
environment's explicit object-at-goal event. REIM learns a nominal policy
$\policyact$, risk predictor $f_\phi$, and recovery policy $\policyrec$. This
semantic-state formulation is the PickPlace mechanism instantiation; the shared
multi-task representation is specified below rather than silently conflated
with the 21-dimensional input.

\subsection{Chunked Nominal Policy}

The state-based ACT model is a conditional variational Transformer that
predicts a chunk of $\ACTChunkLength$ future actions,
\begin{equation}
  \hat a_{t:t+\ACTChunkLength-1}
  = \policyact(s_t,z), \qquad z=0 \text{ at inference}.
\end{equation}
Training minimizes masked L1 reconstruction plus a KL term for the latent
posterior. Overlapping predictions for the current action are combined with
exponentially decaying temporal weights. The ensemble is cleared whenever
control changes between ACT and the recovery option so an unexecuted
pre-intervention chunk cannot re-enter the controller after recovery.

\subsection{Causal Risk Monitoring}

The detector consumes only the most recent $\DetectorHistory$ post-transition
states. For short prefixes, valid states are placed first and the remainder is
zero padded; packed sequence lengths prevent padding from changing the LSTM
state. The failure probability is
\begin{equation}
  p_t = \sigma\!\left(
  g_\phi\!\left(\operatorname{LSTM}_\phi
  (s_{t-h_t+1:t})\right)\right),
  \quad h_t\leq \DetectorHistory .
\end{equation}
Targets indicate whether an online failure event occurs in the inclusive
interval $[t,t+\DetectorForecastHorizon]$. Thus an event observed at the current
step is also positive; causal inputs do not make this a pure future-event
forecast. Binary cross entropy trains the model on disturbed ACT rollouts, and
trajectories rather than individual windows are assigned to disjoint training
and validation splits.

\subsection{Online Trigger-State Recovery Curriculum}
\label{sec:trigger-curriculum}

Approximate failure resets can reproduce an object's displacement without
reproducing the dynamic state that caused it. We instead roll out ACT with the
trained detector under mixed action, observation, and object disturbances. At
the first $p_t\geq\CurriculumTriggerThreshold$, before any recovery action is
executed, the collector stores the complete numeric simulator snapshot:
MuJoCo positions and velocities, mocap targets, actuator/control state, task
goal and reset-time geometry, previous observation, step count, and REIM
episode bookkeeping. Training and model-selection validation use
non-overlapping episode-seed banks. This collection threshold is intentionally
lower than the frozen end-to-end gate
$\tau_{\mathrm{on}}=\FinalControllerTrigger$: the former provides broad
risk-triggered recovery starts, whereas the latter defines the deployed
intervention operating point. We therefore call the states online-triggered,
not deployment-distribution-matched.

A scripted PickPlace expert then continues from the captured state. Every
snapshot is retained for provenance. Only expert-successful trajectories
contribute recovery-imitation targets. This yields
\RecoveryTrainEligibleStarts{} eligible training starts with
\RecoveryTrainExpertSamples{} supervised action pairs and
\RecoveryValidationEligibleStarts{} model-selection starts with
\RecoveryValidationExpertSamples{} pairs. Only 1.12\% of training snapshots
and 1.04\% of validation snapshots record any prior lift. Thus approximately
99\% are pre-grasp or early approach corrections, not post-drop recovery
states.

\subsection{Trigger-Aligned Recovery Imitation}

The recovery option is a deterministic MLP
$\policyrec:\mathbb{R}^{21}\!\rightarrow[-1,1]^4$ with two 256-unit
\texttt{tanh} hidden layers and a clipped linear action head. Let
$\mathcal{D}_{\mathrm{rec}}=\{(s_i,a_i^E)\}$ contain states and scripted-expert
actions from successful continuations after online triggers. We optimize
\begin{equation}
  \mathcal{L}_{\mathrm{rec}}(\psi)=
  \frac{1}{|\mathcal{D}_{\mathrm{rec}}|}
  \sum_{(s,a^E)\in\mathcal{D}_{\mathrm{rec}}}
  \operatorname{SmoothL1}\!\left(\pi_{\mathrm{rec},\psi}(s),a^E\right).
\end{equation}
This is failure-aware, Smooth-L1 behavior cloning: failure awareness enters
through the online LSTM-triggered sampling distribution rather than a
hand-labeled mode bit appended to the actor input.
Training uses 42,386 pairs, 40 epochs, batch size 512, learning rate
$10^{-3}$, and Gaussian observation augmentation with standard deviation
0.005. The epoch with minimum Smooth L1 on 8,212 trajectory-disjoint validation
pairs is retained. Its validation loss is
\RecoveryWarmStartValidationLoss{}.

The deployed file contains only the 21--256--256--4 actor, observation and
action bounds, and provenance metadata. It contains no critic, stochastic
log-standard-deviation, optimizer, or rollout state and records zero
reinforcement-learning environment interactions and zero policy-gradient
updates. An export audit compares the standalone actor against its source on
42,386 training states, 8,212 validation states, and 10,000 deterministic
random-support states: all 60,598 clipped and unclipped action outputs are
bitwise identical.

\subsection{Closed-Loop Arbitration}

At each step the controller updates the causal history and computes $p_t$. ACT
acts while $p_t<\tau_{\mathrm{on}}$; crossing $\tau_{\mathrm{on}}$ activates
the recovery option. The frozen values are
$(\tau_{\mathrm{on}},\tau_{\mathrm{off}},m,B,c)=
(\FinalControllerTrigger,\FinalControllerRelease,
\FinalControllerMinSteps,\FinalControllerBudget,
\FinalControllerClearSteps)$, where $m$ is the minimum commitment, $B$ the
budget, and $c$ the number of clear observations required for release. Because
$m=B$ and $c>B$, the final protocol disables mid-intervention risk-clear
release: $\policyrec$ retains control until task success or budget exhaustion.
If the task remains active after an exhausted intervention, ACT resumes with its
temporal ensemble reset. All evaluation entry points use this same
configuration.

\subsection{Shared Multi-Task Instantiation}

For the breadth study, one task-conditioned controller is learned per benchmark
suite rather than one controller per task. Given the official raw observation
$o_t\in\mathbb{R}^{\MTRawObservationDim}$ and known task index $j$, the shared
input is $x_t=[o_t;e_j]$, where $e_j$ is the official ordered task one-hot. This
gives \MTTenStateDim{} inputs for \MTTenBenchmark{} and \MTFiftyStateDim{} for
\MTFiftyBenchmark{}; the Cartesian action remains four-dimensional. The shared
ACT, causal LSTM, and recovery MLP all receive the same conditioned input.

Multi-task risk supervision is deliberately task agnostic in meaning but is not
identical to the PickPlace event ontology. During data generation only, the
matching Meta-World scripted expert supplies a counterfactual action. For each
task $k$, the behavioral-deviation threshold $\tau_k$ is fitted exclusively on
that task's failure-training bank as the $q=\MTFailureThresholdQuantile$
quantile of ACT--expert action L1 disagreement. The event label is the threshold
crossing, OR the final \MTTerminalFailureTail{} steps of an unsuccessful
episode; causal targets then ask whether an event occurs in the inclusive
window $[t,t+\MTFailureForecastHorizon]$. All $\tau_k$ values are frozen after
this training-bank fit and reused unchanged for validation and final
evaluation. A
fixed raw-disagreement cutoff is retained only as a scale diagnostic, not as a
universal training-label threshold, and the deployed LSTM never accesses expert
actions. At the first detector trigger, the expert takes over in the same online
rollout, and only successful continuations contribute recovery-imitation
targets. Thus the multi-task study preserves learner-induced, trigger-aligned
correction while making no claim that it uses the exact simulator-snapshot
restoration of the PickPlace curriculum.

\begin{figure*}[t]
  \centering
  \PaperGraphic{Figure1_final_framework}{0.96\textwidth}{%
  ACT task policy $\rightarrow$ causal LSTM risk gate $\rightarrow$ supervised
  recovery $\rightarrow$ action execution $\rightarrow$ Meta-World PickPlace
  in MuJoCo (simulated Sawyer model).}
  \caption{REIM training and deployment architecture. The environment block is
  explicitly the Meta-World PickPlace task running in MuJoCo with a simulated
  Sawyer model. Dashed training-only paths denote exact simulator snapshot
  capture, scripted corrective demonstrations, and validation-loss checkpoint
  selection; the scripted expert is never queried during evaluation.}
  \label{fig:framework}
\end{figure*}

\section{Experiments}
\label{sec:experiments}

\subsection{Simulation Environment and Data}

The mechanism study uses the Meta-World 3.1.1 \REIMTask{} single-task,
goal-observable environment~\cite{yu2020metaworld} with MuJoCo
dynamics~\cite{todorov2012mujoco}. The requested benchmark is the historical
PickPlace-v2 Sawyer task. Meta-World 3.1.1 exposes its maintained Gymnasium
successor as \texttt{pick-place-v3}; the wrapper records both identifiers and
maps a v2 request to that maintained registration while preserving the Sawyer
pick--move--place task semantics. We report the exact executed identifier to
avoid implying use of an unavailable legacy API. The 21-dimensional policy
input includes privileged simulator quantities rather than images or a
physical sensing stack.

The ACT dataset contains \REIMDemoTrajectories{} successful scripted
trajectories and \REIMDemoTransitions{} transitions. Failure-predictor data
come from \DetectorTrainingEpisodes{} disturbed ACT episodes, using Gaussian
action and observation noise plus one-shot simulated object-position
displacements. Recovery
starts are collected at \CurriculumNoiseLevel{} nominal disturbance scaling.
The training and validation banks contain \RecoveryTrainStarts{} and
\RecoveryValidationStarts{} exact trigger snapshots, respectively, with
disjoint seed ranges and checkpoint hashes recorded in their manifests.

\subsection{Training Protocol}

ACT uses hidden width \ACTHiddenDim{}, \ACTAttentionHeads{} attention heads,
\ACTEncoderLayers{} encoder layers, \ACTDecoderLayers{} decoder layers, and a
\ACTLatentDim-dimensional latent. It is trained for \ACTTrainingEpochs{}
epochs. The detector has an LSTM width of \DetectorHiddenDim{} followed by an
\DetectorMLPDim-dimensional MLP. The recovery actor uses two 256-unit
\texttt{tanh} layers and is optimized for 40 supervised epochs with batch size
512, learning rate $10^{-3}$, and input-noise standard deviation 0.005.
Training and validation contain 42,386 and 8,212 trigger-aligned pairs,
respectively. Recovery-policy training uses no environment interaction; the
primary uncertainty estimates therefore come from paired held-out episodes,
not multiple policy-gradient seeds.

\subsection{Baselines and Ablations}

The paired baseline comparison contains:
\emph{ACT}, the nominal policy without intervention;
\emph{ACT + Random Reset}, one fresh-task retry after a heuristic failure event
(the retry uses a new simulator seed);
\emph{ACT + Heuristic Recovery}, the same trigger-aligned imitation actor
activated by a semantic heuristic rather than the LSTM; and \emph{REIM}, which
uses the LSTM gate with that actor. The heuristic fires after
at least 20 steps and observed object interaction when a 12-step goal-distance
window improves by less than $10^{-3}$ or worsens by more than
$5\times10^{-3}$; explicit drop/failure flags also fire it. It therefore uses
simulator-level semantic information. The paper-facing component study uses
ACT, heuristic-gated recovery, and full REIM. The first removes recovery
entirely; the second replaces the learned
gate with simulator-level semantics while retaining the identical recovery
actor.

\subsection{Evaluation and Metrics}

The PickPlace main evaluation uses \PlannedMainEpisodes{} paired CRN episode
specifications per method under \CurriculumNoiseLevel{} mixed disturbance,
spanning seeds 8{,}000{,}042--8{,}001{,}041. Robustness uses
\PlannedRobustnessEpisodes{} paired episodes at each of
\RobustnessLevels{}. Each episode records explicit task success, intervention
count, task completion while the recovery controller is active, and length.
For learned recovery methods, ``recovery rate'' is the number of task
completions reached
during active recovery divided by the number of recovery interventions.
Random Reset instead reports post-reset completion per reset and is labeled
separately; the two intervention outcomes are not compared as identical
quantities. Neither is a detector-risk-clearance proxy. Success intervals are
Wilson intervals; ratio and step-count intervals are obtained by episode
bootstrap. Every nominal attempt in a condition consumes the same serialized
common-random-number (CRN) episode specification: task identity, initial
MuJoCo state, disturbance schedule, and episode seed are hash-checked before
aggregation. The Random Reset baseline alone receives one explicitly labeled
fresh-task retry whose second-attempt specification uses a reserved,
deterministically offset seed. The $\tau_{\mathrm{on}}=0.20$ LSTM gate is
frozen before the primary evaluation. A separate post-freeze sweep on 200
previously unused CRN episodes is reported only as a sensitivity diagnostic;
it does not select or alter the primary gate.

\subsection{Known-Task MT10/MT50 Protocol}

The breadth protocol follows Meta-World's goal-observable multi-task setting:
\MTTenBenchmark{} and \MTFiftyBenchmark{} contain known task families, and the
ordered task identity is supplied to every method. It is therefore a shared
known-task evaluation, not an unseen-task, transfer, or meta-learning claim.
Each task contributes the same number of successful expert demonstrations and
\MTPerTaskEvaluationEpisodes{} final episodes. Clean evaluation uses the
official environment success signal and a \MTMaxEpisodeSteps{}-step horizon.
Separate task-bank seeds and payload hashes isolate demonstration, failure,
recovery, validation, and final-evaluation data.

The comparison contains a shared two-layer MLP behavior-cloning capacity
baseline, shared ACT, ACT with the learned recovery actor activated by a
semantic heuristic, and full REIM. The primary statistic is task-macro success;
we also retain per-task success, worst-quartile task success, intervention
burden, and a task-stratified paired bootstrap against ACT. The zero-disturbance
condition is the official-clean anchor. Nonzero levels in
\MTDisturbanceLevels{} apply only task-universal action and observation noise,
never a PickPlace-specific object teleportation, and are labeled a REIM
robustness extension rather than official Meta-World scores. A fresh-task retry
is appendix-only and does not enter the primary multi-task comparison.

\section{Results}
\label{sec:results}

\subsection{PickPlace Mechanism Results}

\input{Table1_final_baseline.tex}

REIM achieves \FinalREIMSuccess{} task success, compared with
\FinalACTSuccess{} for ACT, \FinalResetSuccess{} for random reset, and
\FinalRLRecoverySuccess{} for semantic-heuristic recovery. Across REIM
interventions, \FinalREIMRecovery{} reach task completion while recovery
retains control. Pairing gives a REIM--ACT gain of \FinalREIMGain{}
(\ConfidenceLevel{} CI \FinalREIMGainCI) and a REIM--Random Reset gain of
\FinalREIMGainVsReset{} (CI \FinalREIMGainVsResetCI). Against the stronger
semantic-heuristic gate, REIM improves success by \FinalREIMDeltaVsRL{} (CI
\FinalREIMDeltaVsRLCI) while intervening in \FinalREIMIntervention{} rather
than \FinalRLRecoveryIntervention{} of episodes (655 versus 731 of 1,000).
This corresponds to \FinalREIMFewerIntervenedEpisodes{} fewer intervened
episodes and a
\FinalREIMInterventionReductionVsRL{} relative burden reduction. Thus the
frozen LSTM gate outperforms this simulator-privileged heuristic in the
primary CRN benchmark; the paired gate audit below tests whether the advantage
persists at matched burden.

From no disturbance to the strongest condition, REIM changes from
\FinalREIMSuccessAtZeroNoise{} to \FinalREIMSuccessAtFortyNoise{}, while ACT
changes from \FinalACTSuccessAtZeroNoise{} to
\FinalACTSuccessAtFortyNoise{}. At 10\%, 20\%, 30\%, and 40\% disturbance,
the paired success improvements are \FinalREIMGainAtTenNoise{},
\FinalREIMGainAtTwentyNoise{}, \FinalREIMGainAtThirtyNoise{}, and
\FinalREIMGainAtFortyNoise{}, respectively.

Figure~\ref{fig:gate-sensitivity} compares gate operating points on the
separate diagnostic bank. At the frozen $\tau_{\mathrm{on}}=0.20$ point,
REIM reaches \GateTauTwentySuccess{} success with
\GateTauTwentyIntervention{} intervention, \GateTauTwentyGainVsHeuristic{}
above the matched heuristic (CI \GateTauTwentyGainVsHeuristicCI,
$p=\GateTauTwentyP$, exact two-sided McNemar/binomial test). At
$\tau_{\mathrm{on}}=0.175$, its intervention rate
\GateTauSeventeenFiveIntervention{} nearly matches the heuristic's
\GateHeuristicIntervention{}: the paired difference is
\GateTauSeventeenFiveInterventionDelta{} (CI
\GateTauSeventeenFiveInterventionDeltaCI). REIM nevertheless reaches
\GateTauSeventeenFiveSuccess{} versus \GateHeuristicSuccess{}, a paired gain of
\GateTauSeventeenFiveGainVsHeuristic{} (CI
\GateTauSeventeenFiveGainVsHeuristicCI, $p=\GateTauSeventeenFiveP$). This
post-freeze burden-matched comparison supports the learned gate's utility but
is not substituted for the primary benchmark or used for threshold selection.

\begin{figure*}[t]
  \centering
  \PaperGraphic{Figure2_final_results}{0.96\textwidth}{%
  Paired task-success comparison, intervention burden, and
  \RobustnessLevels{} disturbance robustness curve.}
  \caption{Final simulated Meta-World results. Left: paired task success for
  the four methods. Center: the fraction of episodes receiving an
  intervention, exposing controller burden when success rates are similar.
  Right: success under increasing mixed action and observation noise plus
  simulated one-shot object-position displacements. Every panel is regenerated
  from final per-episode
  records; no development-run curve is substituted.}
  \label{fig:final-results}
\end{figure*}

\begin{figure}[t]
  \centering
  \PaperGraphic{Figure4_gate_sensitivity}{0.98\columnwidth}{%
  Success--intervention sensitivity across five LSTM thresholds.}
  \caption{Post-freeze gate sensitivity on a separate set of 200 paired
  Meta-World seeds at 20\% disturbance. Each point is labeled by
  $\tau_{\mathrm{on}}$; the green cross is the semantic heuristic evaluated on
  the same serialized CRN bank. The circled 0.20 operating point remains the
  frozen primary gate, while the dashed comparison at 0.175 exposes matched
  intervention burden. This diagnostic was not used for model or threshold
  selection.}
  \label{fig:gate-sensitivity}
\end{figure}

\subsection{Known-Task Multi-Task Breadth}

% This table is intentionally invisible until the independent MT10/MT50
% publication gate is enabled after validating all raw episode records.
\ifREIMMultiTaskResults
\begin{table*}[t]
  \centering
  \caption{Official-clean, goal-observable known-task evaluation. All methods
  receive the same official task one-hot, task payloads, and paired episode
  seeds. Success is averaged equally across tasks; brackets are 95\% within-task
  episode-bootstrap intervals. Differences use a paired, task-stratified
  bootstrap against MT-ACT. Worst quartile is the mean of the lowest
  $\lceil T/4\rceil$ task rates, intervention is the task-macro fraction of
  episodes in which recovery acts, and occupancy is the mean fraction of
  episode steps executed by the recovery policy. Nonzero disturbances are
  excluded and reported only as a non-official REIM robustness extension.}
  \label{tab:multitask-clean}
  \scriptsize
  \setlength{\tabcolsep}{4.0pt}
  \begin{tabular}{llccccc}
    \toprule
    Benchmark & Method & Task macro [95\% CI] $\uparrow$ &
      $\Delta$ vs. MT-ACT [95\% CI] & Worst quartile $\uparrow$ & Interv. &
      Occup. $\downarrow$ \\
    \midrule
    \multirow{4}{*}{\MTTenBenchmark{}}
      & MT-MLP BC & \MTTenMLPCleanSuccess{} \MTTenMLPCleanSuccessCI &
        \MTTenMLPDeltaVsACT{} \MTTenMLPDeltaVsACTCI &
        \MTTenMLPWorstQuartile & \NotApplicable & \NotApplicable \\
      & MT-ACT & \MTTenACTCleanSuccess{} \MTTenACTCleanSuccessCI &
        \NotApplicable & \MTTenACTWorstQuartile & \NotApplicable &
        \NotApplicable \\
      & Heuristic-gated recovery &
        \MTTenHeuristicCleanSuccess{} \MTTenHeuristicCleanSuccessCI &
        \MTTenHeuristicDeltaVsACT{} \MTTenHeuristicDeltaVsACTCI &
        \MTTenHeuristicWorstQuartile & \MTTenHeuristicIntervention &
        \MTTenHeuristicOccupancy \\
      & \textbf{MT-REIM (ours)} &
        \textbf{\MTTenREIMCleanSuccess{} \MTTenREIMCleanSuccessCI} &
        \textbf{\MTTenREIMDeltaVsACT{} \MTTenREIMDeltaVsACTCI} &
        \textbf{\MTTenREIMWorstQuartile} & \MTTenREIMIntervention &
        \MTTenREIMOccupancy \\
    \midrule
    \multirow{4}{*}{\MTFiftyBenchmark{}}
      & MT-MLP BC & \MTFiftyMLPCleanSuccess{} \MTFiftyMLPCleanSuccessCI &
        \MTFiftyMLPDeltaVsACT{} \MTFiftyMLPDeltaVsACTCI &
        \MTFiftyMLPWorstQuartile & \NotApplicable & \NotApplicable \\
      & MT-ACT & \MTFiftyACTCleanSuccess{} \MTFiftyACTCleanSuccessCI &
        \NotApplicable & \MTFiftyACTWorstQuartile & \NotApplicable &
        \NotApplicable \\
      & Heuristic-gated recovery &
        \MTFiftyHeuristicCleanSuccess{} \MTFiftyHeuristicCleanSuccessCI &
        \MTFiftyHeuristicDeltaVsACT{} \MTFiftyHeuristicDeltaVsACTCI &
        \MTFiftyHeuristicWorstQuartile & \MTFiftyHeuristicIntervention &
        \MTFiftyHeuristicOccupancy \\
      & \textbf{MT-REIM (ours)} &
        \textbf{\MTFiftyREIMCleanSuccess{} \MTFiftyREIMCleanSuccessCI} &
        \textbf{\MTFiftyREIMDeltaVsACT{} \MTFiftyREIMDeltaVsACTCI} &
        \textbf{\MTFiftyREIMWorstQuartile} & \MTFiftyREIMIntervention &
        \MTFiftyREIMOccupancy \\
    \bottomrule
  \end{tabular}
\end{table*}
\fi


\ifREIMMultiTaskResults
Table~\ref{tab:multitask-clean} reports the official-clean known-task breadth
evaluation. Each task-macro point estimate is accompanied by a 95\%
within-task episode-bootstrap interval. The paired MT-REIM--MT-ACT task-macro differences are
\MTTenREIMDeltaVsACT{} (\ConfidenceLevel{} CI \MTTenREIMDeltaVsACTCI) on
\MTTenBenchmark{} and \MTFiftyREIMDeltaVsACT{} (CI
\MTFiftyREIMDeltaVsACTCI) on \MTFiftyBenchmark{}. Nonzero-disturbance results
are reported separately as a robustness extension and are not described as
official clean benchmark scores. These noise levels use the same absolute
scales as the single-task protocol ($\sigma_a = 0.4\lambda$, $\sigma_o =
0.025\lambda$ at level $\lambda$) but omit object-position displacements;
MT-ACT's sharp degradation at $\lambda = 0.1$ therefore reflects the shared
task-conditioned policy's sensitivity rather than a scale mismatch.

\begin{figure*}[t]
  \centering
  \PaperGraphic{Figure_multitask_robustness}{0.94\textwidth}{%
  MT10 and MT50 task-macro success under task-universal disturbance.}
  \caption{Non-official REIM robustness extension on the known-task,
  goal-observable MT10 and MT50 suites. Curves apply task-universal Gaussian
  action and observation noise; they do not alter object poses. Points are
  task-macro success over 50 episodes per task, and bars are 95\%
  within-task episode-bootstrap intervals. These disturbed conditions are
  deliberately separated from the official-clean scores in
  Table~\ref{tab:multitask-clean}.}
  \label{fig:multitask-robustness}
\end{figure*}
\else
\paragraph{Result status.}
The MT10/MT50 protocol and artifact schema are implemented, but complete clean
and disturbed episode records have not yet passed the independent publication
gate. Consequently, this manuscript currently makes no MT10/MT50 success,
intervention, comparative, or robustness-trend claim, and the numerical
multi-task table is intentionally suppressed.
\fi

\subsection{Completed Component Diagnostics}

\begin{table}[t]
  \centering
  \caption{Completed component diagnostics. These values use different
  validation units and are not end-to-end comparisons. The recovery row uses
  trajectory-disjoint exact-trigger validation pairs and is therefore
  model-selection evidence rather than an unbiased test result.}
  \label{tab:component-diagnostics}
  \small
  \setlength{\tabcolsep}{4pt}
  \begin{tabular}{lll}
    \toprule
    Component & Diagnostic & Value \\
    \midrule
    ACT & validation action L1 & \ACTBestValidationLone \\
    LSTM & precision at $\tau=\DetectorDiagnosticThreshold$ &
      \DetectorPrecisionDiagnostic \\
    LSTM & recall at $\tau=\DetectorDiagnosticThreshold$ &
      \DetectorRecallDiagnostic \\
    Recovery imitation & validation Smooth L1 &
      \RecoveryWarmStartValidationLoss \\
    Standalone actor & audited output error & exactly zero \\
    \bottomrule
  \end{tabular}
\end{table}


ACT reaches validation action L1 \ACTBestValidationLone{} on a
trajectory-grouped split. At the separately reported
\DetectorDiagnosticThreshold{} detector operating point, precision is
\DetectorPrecisionDiagnostic{} and recall is \DetectorRecallDiagnostic{} over
\DetectorValidationSamples{} validation windows. Figure~\ref{fig:detector}
shows the corresponding confusion matrix. The diagnostic is evaluated at the
same frozen probability threshold as end-to-end arbitration,
$\tau_{\mathrm{on}}=\FinalControllerTrigger$, but window-level classification
metrics do not by themselves establish closed-loop utility.

The inclusive target requires a second, stricter audit. Of the positive
validation windows, \DetectorCurrentEventPositiveShare{} have event offset zero.
After excluding those current-event windows and scoring only future offsets
$1$--$\DetectorForecastHorizon$ against windows with no event in the horizon,
precision/recall/F1 are \DetectorStrictPreEventPrecision{}/
\DetectorStrictPreEventRecall{}/\DetectorStrictPreEventFone{}. At trajectory
level, the monitor alerts strictly before the first event in
\DetectorPreEventDetectedTrajectories{}/\DetectorEventTrajectories{}
(\DetectorPreEventTrajectoryDetection{}) event trajectories, with median latest
alert lead \DetectorMedianPreEventLead{} step. We therefore describe the LSTM as
a causal risk monitor, not a calibrated long-horizon forecaster.

\begin{figure}[t]
  \centering
  \PaperGraphic{Figure3_detector}{0.86\columnwidth}{%
  Grouped-validation confusion matrix.}
  \caption{Causal risk-monitor confusion matrix on the trajectory-grouped
  validation split at $\tau=\DetectorDiagnosticThreshold$. Counts are
  $(\mathrm{TN},\mathrm{FP},\mathrm{FN},\mathrm{TP})=
  (\DetectorTNDiagnostic,\DetectorFPDiagnostic,
  \DetectorFNDiagnostic,\DetectorTPDiagnostic)$.}
  \label{fig:detector}
\end{figure}

The recovery actor reaches validation Smooth L1
\RecoveryWarmStartValidationLoss{}. The standalone checkpoint records zero
environment-training steps and contains only 72,452 actor parameters. Its
export audit finds exactly zero action error on 60,598 states, including every
training and validation state plus 10,000 random-support states. This
establishes deployment-artifact equivalence; held-out paired episodes, rather
than the export check, support the task-level claims.

\begin{figure*}[t]
  \centering
  \PaperGraphic{Figure5_operation_sequence}{0.96\textwidth}{%
  Qualitative sequence from one separate Meta-World/MuJoCo episode:
  ACT execution, LSTM trigger, recovery-imitation intervention, and task
  completion.}
  \caption{Paired qualitative operation sequence from separate simulation seed
  \OperationTraceSeed{} at \OperationTraceNoise{} disturbance. With the
  frozen evaluation gate
  $\tau_{\mathrm{on}}=\OperationTraceThreshold$, ACT does not complete the
  task within \OperationTraceACTSteps{} steps, whereas REIM triggers once and
  completes in \OperationTraceREIMSteps{} steps, including
  \OperationTraceRecoverySteps{} supervised-recovery actions. These are rendered
  Meta-World/MuJoCo frames of the simulated Sawyer model, not hardware
  photographs. This selected trace explains controller operation and does not
  estimate aggregate success.}
  \label{fig:operation}
\end{figure*}

\section{Recovery and Gate Ablations}
\label{sec:ablations}

\input{Table2_final_ablation.tex}

Removing recovery entirely reduces the controller to ACT and
\AblationASuccess{} success. Keeping the same trigger-aligned recovery actor
but replacing the LSTM with simulator-semantic heuristics reaches
\AblationCSuccess{} and intervenes in \FinalRLRecoveryIntervention{} of
episodes. Full REIM reaches \AblationDSuccess{} while intervening in only
\FinalREIMIntervention{}, improving both success and intervention burden. The
separate burden-matched audit in Fig.~\ref{fig:gate-sensitivity} makes the gate
comparison under nearly equal intervention rates.

\begin{figure}[t]
  \centering
  \PaperGraphic{Figure3_final_ablation}{0.98\columnwidth}{%
  Recovery and gate ablation on paired simulated episodes.}
  \caption{Paper-facing component removal study. ACT removes recovery;
  Heuristic gate + recovery replaces only the learned LSTM arbitration rule
  while retaining the same recovery actor; REIM uses LSTM-triggered recovery
  imitation. Error bars are Wilson 95\% intervals.}
  \label{fig:ablation}
\end{figure}

Table~\ref{tab:final-ablation} therefore supports two scoped conclusions:
recovery is responsible for the large improvement over nominal ACT, and learned
temporal arbitration improves on a strong simulator-privileged gate. It does
not isolate online snapshot collection from recovery-imitation training, so we
make no separate marginal claim about those coupled design choices.

\section{Limitations}
\label{sec:limitations}

The current study has several boundaries. First, its mechanism evidence and
component ablations concentrate on state-observed PickPlace. The MT10/MT50
breadth protocol supplies known task identity and therefore cannot establish
unseen-task generalization, meta-learning, or transfer; visual perception and
contact sensing also remain outside scope. Second, exact simulator restoration
in the PickPlace curriculum depends on privileged MuJoCo state and task
bookkeeping. It is a controlled way to align recovery training in simulation,
not a deployable physical-robot reset mechanism. Hardware would require state
estimation and safe physical repositioning or real failure collection. Third,
filtering recovery-imitation
trajectories by scripted-expert success biases training toward recoverable states and may
exclude precisely the hardest failures. Fourth, detector metrics are measured
on rollouts from a fixed nominal policy in each protocol and may shift as ACT or
the disturbance distribution changes. Moreover, the PickPlace inclusive label
makes the standard confusion matrix predominantly a current-event diagnostic, and
\DetectorNonEventTrajectoryAlert{} of validation trajectories without a labeled
event still contain an alert at the deployed threshold. Fifth, the main table
is conditional on one trained ACT, detector, and supervised recovery actor; the
multi-task protocol likewise uses one trained stack per suite. Paired episode
uncertainty therefore does not propagate all model-training uncertainty. Sixth, the validation
pairs select the recovery epoch and therefore provide model-selection rather
than independent test evidence. Finally, the reported ``intervention outcome''
is task completion while the recovery option remains in control: the
persistent option does not demonstrate recovery
to a certified safe set followed by an ACT hand-back. It can also waste steps
after false positives, so all arbitration parameters remain frozen throughout
the paired benchmark.

\section{Conclusion}
\label{sec:conclusion}

REIM combines state-based chunked imitation, causal runtime risk monitoring,
and a persistent trigger-aligned imitation recovery option in one closed-loop
manipulation controller. Corrective demonstrations are collected from complete
simulator states captured at online risk-monitor triggers, and a deterministic
recovery actor is selected by trajectory-disjoint validation loss without
reinforcement-learning interaction. In paired PickPlace Meta-World/MuJoCo
evaluation, REIM improves ACT by \FinalREIMGain{} and remains
substantially more robust at strong disturbances. Its LSTM gate reduces
interventions by \FinalREIMInterventionReductionVsRL{} relative to
heuristic-triggered recovery while improving success by
\FinalREIMDeltaVsRL{} in the primary paired benchmark; a separate
burden-matched audit yields \GateTauSeventeenFiveGainVsHeuristic{} at
$p=\GateTauSeventeenFiveP$.
A shared task-conditioned instantiation additionally defines a known-task
MT10/MT50 breadth test whose official-clean evidence is kept separate from the
task-universal disturbance extension.
On the clean known-task conditions the learned gate is statistically level
with the semantic heuristic and the heuristic is slightly ahead, so the gate's
advantage appears at higher disturbance rather than as a universal improvement
over the heuristic.
The captured distribution is approximately 99\% pre-grasp/early correction,
so the evidence supports selective simulated failure prevention and recovery,
not general post-drop repair, visual perception, sim-to-real transfer, or a
physical Sawyer claim.

\appendix

\section{Multi-Seed Stability of the Supervised Recovery Gate}
\label{app:multiseed}

To probe whether the multi-task findings hinge on one lucky training seed, we
trained the supervised MT10 recovery gate on three seeds (42, 43, 44) and
evaluated each on the same MT10 bank (20265010, 50 episodes per task).
Table~\ref{tab:mt10-seed-stability} and Fig.~\ref{fig:mt10-seed-stability}
show that seed-to-seed variation is small relative to the disturbance
conditioning itself: under clean conditions the supervised gate tracks the
heuristic within one point, while at 40\% disturbance its advantage over the
heuristic gate is consistent across seeds. The paired-vs-ACT rescued/harmed
counts (mean $\pm$ s.d. over the three seeds) corroborate this: at 40\%
disturbance REIM rescues 282.3 $\pm$ 24.5 episodes and harms 1.3 $\pm$ 1.2,
versus the heuristic gate rescuing 181.3 $\pm$ 17.9 and harming 0.3 $\pm$ 0.6.
This analysis concerns the current supervised recovery gate; the earlier PPO
recovery negative control remains excluded from this paper.

\begin{table}[t]
  \centering
  \caption{MT10 multi-seed stability of the supervised recovery gate
  (bank 20265010, 50 episodes per task per seed). Rescued/Harmed count
  episodes where the method flips the paired MT-ACT outcome.
  Seed-to-seed variation is small relative to the noise-0.4 gap over
  the heuristic gate.}
  \label{tab:mt10-seed-stability}
  \small
  \setlength{\tabcolsep}{3.4pt}
  \begin{tabular}{llcccc}
    \toprule
    Noise & Method & Seed 42 & Seed 43 & Seed 44 & Mean $\pm$ s.d. \\
    \midrule
    clean & MT-REIM & 97.0\% & 96.6\% & 96.2\% & 96.6 $\pm$ 0.4\% \\
    clean & Heuristic gate & 96.6\% & 95.8\% & 97.6\% & 96.7 $\pm$ 0.9\% \\
    \midrule
    0.1 & MT-REIM & 45.6\% & 43.0\% & 43.8\% & 44.1 $\pm$ 1.3\% \\
    0.1 & Heuristic gate & 47.8\% & 41.4\% & 45.0\% & 44.7 $\pm$ 3.2\% \\
    \midrule
    0.4 & MT-REIM & 61.4\% & 55.6\% & 65.4\% & 60.8 $\pm$ 4.9\% \\
    0.4 & Heuristic gate & 45.0\% & 38.8\% & 38.6\% & 40.8 $\pm$ 3.6\% \\
    \bottomrule
  \end{tabular}
\end{table}


\begin{figure}[t]
  \centering
  \PaperGraphic{FigureA1_mt10_seed_stability}{0.98\columnwidth}{%
  MT10 seed stability of the supervised recovery gate.}
  \caption{MT10 multi-seed stability of the supervised recovery gate.
  Bars give the task-macro success mean over three training seeds
  (42/43/44) with $\pm 1$ s.d. whiskers; dots are the per-seed values.
  Evaluation uses bank 20265010 with 50 episodes per task per seed.}
  \label{fig:mt10-seed-stability}
\end{figure}

\section{Recovery Training-Design Controls}
\label{app:training-design}

Table~\ref{tab:training-design} and Fig.~\ref{fig:training-design} audit five
training-design choices for the recovery actor on 200 paired Meta-World
development seeds at 20\% disturbance, contrasting the supervised REIM
imitation baseline with PPO-based controls. Exact online trigger-state
resets, expert actor warm start, and task-selected checkpoints each matter;
removing the warm start collapses success, and an unselected 500k endpoint
underperforms the selected checkpoint. These controls justify the training
design but use a development protocol (200 paired episodes, one training
seed), so they are reported here rather than in the main benchmark claims.

\begin{table}[t]
  \centering
  \caption{Recovery training-design controls on the same 200 paired
  Meta-World seeds at 20\% disturbance. The REIM actor receives zero RL
  interactions. PPO controls receive at least 500k environment steps;
  ``selected'' denotes a development-bank checkpoint and ``endpoint'' the
  actual 500k endpoint.}
  \label{tab:training-design}
  \scriptsize
  \setlength{\tabcolsep}{3.2pt}
  \begin{tabular}{lcccc}
    \toprule
    Condition & Start & Init. & Training & Success $\uparrow$ \\
    \midrule
    REIM imitation (ours) & Online & Expert & Supervised & 91.5\% \\
    PPO selected control & Online & Expert & PPO selected & 91.0\% \\
    PPO approximate starts & Approx. & Expert & PPO selected & 89.5\% \\
    PPO without imitation & Online & Random & PPO selected & 54.5\% \\
    PPO endpoint & Online & Expert & PPO endpoint & 75.5\% \\
    \bottomrule
  \end{tabular}
\end{table}


\begin{figure*}[t]
  \centering
  \PaperGraphic{Figure7_training_design}{0.98\textwidth}{%
  Training-design ablation: success, recovery behavior, and paired gain.}
  \caption{Recovery training-design controls on the same 200 paired
  Meta-World seeds at 20\% disturbance. (a) End-to-end success with 95\%
  Wilson intervals; the dashed line marks the ACT reference. (b) Recovery
  intervention outcome and intervened-episode rates. (c) Paired success
  difference over ACT with 95\% bootstrap intervals. PPO controls receive at
  least 500k environment steps; the actor-only row receives zero.}
  \label{fig:training-design}
\end{figure*}

\section{Per-Episode Recovery Timelines}
\label{app:recovery-timelines}

Fig.~\ref{fig:recovery-timelines} shows representative per-episode timelines
from the formal CRN evaluation bank (seeds from 8000042, 20\% disturbance),
complementing the qualitative montage of Fig.~\ref{fig:operation} with quantitative traces of
detector score, gate state, and recovery execution. Successful episodes show
a triggered gate followed by a persistent recovery option that restores
progress; failed episodes exhaust the recovery budget without re-grasp.

\begin{figure*}[t]
  \centering
  \PaperGraphic{Figure4_recovery}{0.96\textwidth}{%
  Per-episode detector and recovery timelines.}
  \caption{Per-episode recovery timelines from the formal CRN bank
  (20\% disturbance). Each column is one episode: the top row tracks the
  detector risk score and gate events, and the bottom row shows recovery
  execution. Episode captions report the recovery outcome.}
  \label{fig:recovery-timelines}
\end{figure*}

\balance
\bibliographystyle{IEEEtran}
\bibliography{reim_refs}

\end{document}
